% !TeX root = ../../Book2.tex
% 纲要标题：## 第11章　代数、反环与双中心化子
% 纲要位置：计划纲要.md，第 1406 行。

\chapter{代数、反环与双中心化子}
\label{b2:ch:11}
\BookChapterNumber{b2:ch:11}

\begin{chapterabstract}
代数同时具有线性结构与乘法，表示则把乘法实现为线性算子的复合。本章通过结构常数、群代数与矩阵代数说明这两种结构如何相容，再用反环把右作用写成左作用，用包络代数把双模写成左模。幂等元给出正则模的分解和角环中的乘法；中心化子由与给定作用交换的算子组成。通过正则作用和张量因子的显式计算，建立双中心化子的基本例子，并辨明其何时不能恢复原作用代数。
\end{chapterabstract}

\begin{learninggoals}
\begin{itemize}
\item 根据结构常数检查结合律和单位元，计算换基后的乘法，区分向量空间同构与代数同构。
\item 在群表示、群代数模、双模及包络代数模之间转换，并核对反环与复合的次序。
\item 从正交幂等元构造模分解和矩阵分块，解释中心性何时把模分解提升为代数乘积分解。
\item 计算给定作用的中心化子及双中心化子，区分交换作用、互为中心化子与不可约性。
\end{itemize}
\end{learninggoals}

前两编研究了模的生成、关系与分解，现在要进一步问：一个代数可以怎样作用在这些模上，哪些模自同态与作用交换，又能否从这些自同态恢复原代数？把矩阵代数 \(A=M_2(k)\) 看作自身的向量空间，左乘 \(L_a(x)=ax\) 与右乘 \(R_b(x)=xb\) 总是交换：\(a(xb)=(ax)b\)。不过右乘的复合顺序反过来，\(R_b\circ R_c=R_{cb}\)：先作用 \(R_c\)，再作用 \(R_b\)。若要把两侧作用同时写成代数作用，就须用反环记录右侧的乘法方向。

\ProseParagraphBreak

同一组矩阵还可以从另一面观察。幂等矩阵把正则模分成像与核，和所有左乘都交换的算子则组成右作用的中心化子。先在矩阵单位上逐项计算这些事实，再抽去坐标，才知道哪些由代数结构保证、哪些只是在这个特殊例子里成立。

\ProseParagraphBreak

本章使用\chapref{b2:ch:02}的模与自同态环、\chapref{b2:ch:06}的双模和张量积，以及\chapref{b2:ch:08}的结合代数与张量作用。底域始终写作 \(k\)，代数同态保持单位元；一般不假设 \(k\) 代数闭，也不默认代数交换或有限维。矩阵单位提供了可直接核验的计算方法；左乘、右乘与模自同态产生的投影需要分别辨认。

\ProseParagraphBreak
代数和表示的基本语言可参阅 Etingof 等人的《Introduction to Representation Theory》\cite[\S 2.2, pp. 8--9; \S 2.3, pp. 9--10]{EtingofEtAl2011RepresentationTheory}。该书在这一部分默认底域代数闭并排除零代数；这里允许任意底域和零代数。后面的半单性与稠密性理论将进一步解释双中心化子计算。

% !TeX root = ../../Book2.tex
% 纲要标题：### 11.1 域上的代数
% 纲要位置：计划纲要.md，第 1408 行。

\section{域上的代数}
\label{b2:sec:1101}

矩阵可以相加、数乘，也可以相乘；群元素的形式线性组合也具有这三种运算。两者都是代数：乘法与向量空间的加法、数乘相容。本章以域 \(k\) 为底域，除特别说明外，代数结合且含单位元，同态保持单位元；允许零代数，需要把底域嵌入代数或讨论非零简单对象时，再明确加上非零条件。

% 纲要要点（供目录核对）：
%
% * 中心标量、结合含幺代数及表示与模同态
% * 乘法映射的结构常数、结合方程及基变换
% * 群代数泛性质、增广与矩阵单位作用
%

\subsection{结合含幺代数与代数同态}
\label{b2:subsec:1101:01}

\begin{definition}\label{b2:def:field-algebra}
设 \(k\) 为域。若 \(k\) 向量空间 \(A\) 配有双线性乘法 \((a,b)\mapsto ab\) 和元素 \(1_A\)，且对任意 \(a,b,c\in A\) 满足
\[
(ab)c=a(bc),\qquad 1_Aa=a1_A=a.
\]
则称 \(A\) 为域 \(k\) 上的结合含幺代数。若 \(A,B\) 均为这样的代数，\(k\) 线性映射 \(f:A\to B\) 保持乘法与单位元，则称 \(f\) 为 \(k\) 代数同态。\(A\) 中与所有元素交换的元素组成其中心 \(Z(A)\)。
\end{definition}

这是\cref{b2:def:associative-graded-algebra}中交换底环代数的定义在底环为域 \(k\) 时的具体形式。双线性给出 \((\lambda1_A)a=\lambda a=a(\lambda1_A)\)，所以 \(\lambda\mapsto\lambda1_A\) 是到中心的保单位元环同态；反过来，这样的结构同态给出标量乘法及上述双线性。环结构与线性结构正是通过中心的标量元素相容的。若 \(A\ne0\)，这个同态单射，因为其核是域的真理想。以下直接把 \(\lambda\in k\) 视作 \(A\) 中的元素时，默认 \(A\ne0\)；零代数仍由结构同态解释标量作用。

\ProseParagraphBreak
例如 \(k[t]/(t^2)\) 的元素唯一写成 \(a+b\varepsilon\)，其中 \(\varepsilon^2=0\)。它作为向量空间与 \(k\times k\) 都是二维，但前者有非零幂零元，后者没有。代数同构保持幂零性，所以它们不同构。区分代数应当寻找这样的乘法性质，不能只比较向量空间维数或某组基下的乘法表。

\begin{definition}\label{b2:def:algebra-representation}
设 \(k\) 为域、\(A\) 为含幺结合 \(k\) 代数、\(V\) 为 \(k\) 向量空间。保单位元的 \(k\) 代数同态 \(\rho:A\to\operatorname{End}_k(V)\) 称为 \(A\) 在 \(V\) 上的表示。若 \(V\ne0\)，且 \(V\) 没有非零真 \(A\) 稳定子空间，就称该表示为\term[bukeyue-biaoshi]{不可约表示}[irreducible representation]；此时相应左 \(A\) 模称为\term[jiandan-mo]{简单模}[simple module]。
\end{definition}

本章用不可约表示和简单模的定义辨认基本作用；它们与半单模的结构将在\chapref{b2:ch:13}讨论。

\begin{proposition}\label{b2:prop:algebra-modules-representations}
设 \(k\) 为域，\(A\) 为含幺结合 \(k\) 代数，\(V\) 为 \(k\) 向量空间。在 \(V\) 上指定幺性左 \(A\) 模结构，并使底域作用与已有标量作用相同，等价于指定保单位元的 \(k\) 代数同态
\[
\rho:A\longrightarrow\operatorname{End}_k(V).
\]
对两个这样的左 \(A\) 模 \(V,W\)，将相应表示记为 \(\rho_V,\rho_W\)。模同态正是满足 \(f\rho_V(a)=\rho_W(a)f\)（\(a\in A\)）的 \(k\) 线性映射 \(f:V\to W\)。
\end{proposition}
\begin{proof}
由模结构令 \(\rho(a)(v)=av\)。分配律给出对 \(a\) 和 \(v\) 的可加性；底域在中心保证每个 \(\rho(a)\) 为 \(k\) 线性。模的结合律与幺性给出 \(\rho(ab)=\rho(a)\rho(b)\)、\(\rho(1)=\operatorname{id}\)，而 \(\rho(\lambda a)=\lambda\cdot\rho(a)\)。反向用同一公式定义作用，这些等式逐项给出模公理。最后的相容等式在 \(v\) 处就是 \(f(av)=a\cdot f(v)\)。
\end{proof}

一个子空间是 \(A\) 子模，正是说它在每个 \(\rho(a)\) 下稳定。令上面的 \(W=V\)，模自同态的条件就成为
\[
\operatorname{End}_A(V)
=\{\,T\in\operatorname{End}_k(V)\mid T\rho(a)=\rho(a)T\;\text{对所有}\;a\in A\,\}.
\]
它由与全部作用算子交换的算子组成，\secref{b2:sec:1104}将把这种集合称为中心化子。表示的核为 \(V\) 的零化子；核为零恰好表示作用忠实。一个表示可以使许多不同代数元素产生同一个线性算子，定义本身并不排除这种情况。

\subsection{结构常数及基变换}
\label{b2:subsec:1101:02}

设 \(A\) 为域 \(k\) 上的有限维结合代数，以 \(e_1,\ldots,e_n\) 为一组 \(k\) 基。乘法的双线性与\cref{b2:thm:tensor-universal}给出 \(k\) 线性映射 \(\mu:A\otimes_kA\to A\)、\(\mu(a\otimes b)=ab\)。逐对展开乘积，写成
\[
\mu(e_i\otimes e_j)=e_ie_j=\sum_{r=1}^n c_{ij}^{r}e_r.
\]
所得标量 \(c_{ij}^{r}\in k\) 称为 \(A\) 在这组基下的\term[jiegou-changshu]{结构常数}[structure constants]。两个下标记录乘法的两个输入，上标记录输出坐标。它们决定全部乘法，因为两个一般元素的乘积可由双线性展开。

\begin{proposition}\label{b2:prop:structure-constants}
设 \(V\) 是域 \(k\) 上以 \(e_1,\ldots,e_n\) 为基的向量空间。指定 \(k\) 中的常数 \(c_{ij}^{r}\)（\(1\le i,j,r\le n\)），并规定 \(e_ie_j=\sum_r c_{ij}^{r}e_r\)，就指定了 \(V\) 上唯一的双线性乘法。该乘法结合，当且仅当对所有 \(i,j,\ell,s\in\{1,\ldots,n\}\)，
\begin{equation}\label{b2:eq:structure-associativity}
\sum_r c_{ij}^{r}c_{r\ell}^{s}=\sum_r c_{j\ell}^{r}c_{ir}^{s}.
\end{equation}
元素 \(u=\sum_i u_i e_i\) 是单位元，当且仅当
\[
\sum_i u_i c_{ij}^{s}=\delta_{js},\qquad
\sum_i u_i c_{ji}^{s}=\delta_{js}
\quad\text{对所有}\;j,s.
\]
\end{proposition}
\begin{proof}
令两个线性组合的乘积为其所有基乘积按系数展开的有限和。坐标唯一使其良定义，并直接得到双线性。比较 \((e_ie_j)e_\ell\) 与 \(e_i(e_je_\ell)\) 的第 \(s\) 个坐标，便得到\eqref{b2:eq:structure-associativity}。两种三变量表达都三线性，故在所有基向量组上相等即在任意输入上相等。类似地，\(ue_j=e_j\) 与 \(e_ju=e_j\) 的坐标恰是两条单位条件；线性性再推广到任意元素。
\end{proof}

因此，固定有限维向量空间上的结合乘法由满足\eqref{b2:eq:structure-associativity}的结构常数描述；这些是关于 \(c_{ij}^{r}\) 的二次多项式方程。乘法确定以后，寻找单位元只需解关于 \(u_i\) 的线性方程；若同时把 \(u_i\) 与 \(c_{ij}^{r}\) 当作未知量，单位条件则是关于这两组变量的双线性关系。

\ProseParagraphBreak
若换基为 \(f_a=\sum_i p_{ia}e_i\)，记 \(P^{-1}=(q_{dr})\)，两个输入向量的坐标都先用 \(P\) 转到旧基，乘积的输出坐标再用 \(P^{-1}\) 转回新基。因此新结构常数为
\begin{equation}\label{b2:eq:structure-basis-change}
(c')_{ab}^{d}=\sum_{i,j,r}q_{dr}p_{ia}p_{jb}c_{ij}^{r}.
\end{equation}
具体地，将 \(f_a,f_b\) 分别展开后，得到 \(f_af_b=\sum_{i,j,r}p_{ia}p_{jb}c_{ij}^{r}e_r\)，再代入 \(e_r=\sum_dq_{dr}f_d\) 即得此式。两个 \(P\) 和一个 \(P^{-1}\) 分别对应两个输入与一个输出，不能按只有一个输入的算子来作相似变换。

\ProseParagraphBreak
例如，在双数代数中以 \(1,y\) 为新基，令 \(y=b+a\varepsilon\)、\(a\ne0\)。直接计算得到 \(y^2=2by-b^2\)。乘法表中虽然出现了常数项和一次项，元素 \(y-b=a\varepsilon\) 仍非零且平方为零。因此结构常数随基改变，而存在非零幂零元这一乘法性质不变。

\ProseParagraphBreak
结构常数回答固定有限维向量空间上可以怎样定义乘法。群代数和矩阵代数则已经给定乘法；研究它们在向量空间上的作用，要检验乘法是否对应算子复合，并找出作用下保持不变的子空间。

\subsection{群代数和矩阵代数}
\label{b2:subsec:1101:03}

群乘法只规定群元素之间怎样相乘。以这些元素为基，允许有限线性组合，再将乘法双线性延拓，就把群乘法变成代数乘法，同时保留了原来的每个群元素。

\begin{definition}\label{b2:def:group-algebra}
设 \(k\) 为域、\(G\) 为群。以 \(G\) 的元素为基构造 \(k\) 向量空间，其元素为有限和 \(\sum_g a_g g\)（\(a_g\in k\)）。把群乘法按 \(k\) 双线性延伸：
\[
\left(\sum_g a_g g\right)\left(\sum_h b_h h\right)
=\sum_{g,h}a_gb_h(gh).
\]
所得含幺 \(k\) 代数称为群 \(G\) 的\term[qun-daishu]{群代数}[group algebra]，记为 \(k[G]\)；其单位为群单位元对应的基向量。
\end{definition}
\symidx{\(k[G]\)：群 \(G\) 的群代数}

两个支集有限，乘积仍是有限和；三重乘积的系数由同一个有限三重和给出，群乘法的结合律保证结果相同。因此该式定义结合代数。它交换当且仅当群交换，因为不同群元素是不同的基向量。若 \(G\) 有限，则 \(\dim_k k[G]=|G|\)。

\begin{proposition}\label{b2:prop:group-algebra-representations}
设 \(k\) 为域、\(G\) 为群、\(B\) 为非零结合含幺 \(k\) 代数，\(B^\times\) 为其单位群。限制到基向量 \(g\in G\) 给出双射
\[
\operatorname{Hom}_{k\text{-alg}}(k[G],B)
\longrightarrow\operatorname{Hom}_{\mathrm{Grp}}(G,B^\times).
\]
其逆把群同态 \(\rho:G\to B^\times\) 线性延拓为
\[
\widetilde\rho\left(\sum_g a_g g\right)=\sum_g a_g\cdot\rho(g).
\]
此对应对 \(G\) 逆变、对 \(B\) 协变地自然：对群同态 \(\alpha:H\to G\) 及非零结合含幺 \(k\) 代数之间的同态 \(\psi:B\to C\)，预复合 \(k[\alpha]:k[H]\to k[G]\) 对应预复合 \(\alpha\)，后复合 \(\psi\) 对应后复合其单位群限制 \(\psi^\times:B^\times\to C^\times\)。这里 \(k[\alpha]\) 是 \(h\mapsto\alpha(h)\) 的线性延拓。

对任意 \(k\) 向量空间 \(V\)（允许 \(V=0\)），群表示 \(G\to\operatorname{GL}(V)\) 与 \(V\) 上底域作用为原标量作用的幺性左 \(k[G]\) 模结构一一对应；两个表示之间的等变 \(k\) 线性映射对应 \(k[G]\) 模同态。
\end{proposition}
\begin{proof}
对代数同态 \(f:k[G]\to B\)，群基向量在群代数中满足 \(gg^{-1}=g^{-1}g=1\)，故 \(f(g)f(g^{-1})=f(g^{-1})f(g)=1_B\)。所以 \(f(g)\) 确在 \(B^\times\) 中，且保持群乘法。反向给定 \(\rho:G\to B^\times\)，由于 \(G\) 是一组基，其线性延拓 \(\widetilde\rho\) 唯一。由有限和展开，
\[
\begin{aligned}
\widetilde\rho\left(\left(\sum_g a_g g\right)\left(\sum_h b_h h\right)\right)
&=\sum_{g,h}a_gb_h\cdot\rho(gh)\\
&=\left(\sum_g a_g\cdot\rho(g)\right)\left(\sum_h b_h\cdot\rho(h)\right).
\end{aligned}
\]
因此它保持乘法；又 \(\rho(1_G)=1_B\)，故它是保单位元的代数同态。限制与线性延拓互逆，因为线性映射在一组基上的取值决定整个映射。

群同态 \(\alpha\) 的线性延拓同样保持乘法与单位元，因而 \(k[\alpha]\) 是代数同态；\(\psi\) 把单位及其逆仍送到互逆元素，所以单位群限制 \(\psi^\times\) 定义良好。对每个 \(h\in H\)，复合的取值都为 \(\psi(\rho(\alpha(h)))\)。由基上的唯一性，
\[
\widetilde{\psi^\times\circ\rho\circ\alpha}
=\psi\circ\widetilde\rho\circ k[\alpha],
\]
这就证明了两个变量的自然性。

当 \(V\ne0\) 时，取 \(B=\operatorname{End}_k(V)\)，其单位群就是 \(\operatorname{GL}(V)\)，再由\cref{b2:prop:algebra-modules-representations}得到表示与模结构的对应。明确地，作用为 \((\sum_g a_g g)v=\sum_g a_g\cdot\rho(g)v\)；反向限制到 \(g\) 的作用算子，其逆为 \(g^{-1}\) 的作用，因为 \(g\cdot(g^{-1}\cdot v)=v=g^{-1}\cdot(g\cdot v)\)。若 \(V=0\)，\(\operatorname{GL}(0)\) 是单元素群，且零空间只有一个幺性左模结构，对应仍成立。对两个空间之间的 \(k\) 线性映射，与每个 \(g\) 的作用相容等价于与它们的任意有限线性组合的作用相容，故等变映射与模同态也一致。
\end{proof}

把每个群元素都送到 \(1\in k^\times\)，就忘掉了群元素之间的区别。对群 \(G\) 和域 \(k\)，由上述泛性质得到的代数同态 \(\epsilon:k[G]\to k\)、\(\sum_g a_g g\mapsto\sum_g a_g\) 只保留系数总和，称为\term[zengguang]{增广}[augmentation]。相应一维模上每个群元素都恒等地作用，正是平凡表示。其核由 \(g-1\)（\(g\in G\)）线性生成，因为系数和为零的元素可写成 \(\sum_g a_g(g-1)\)。

\ProseParagraphBreak
矩阵代数 \(M_n(k)\)、\(n\ge1\)，以矩阵单位 \(E_{ij}\) 为基，乘法满足
\[
E_{ij}E_{rs}=\delta_{jr}E_{is},\qquad 1=\sum_iE_{ii}.
\]
\symidx{\(M_n(k)\)：域上的 \(n\) 阶矩阵代数}
它在列向量空间 \(k^n\) 上自然作用。矩阵单位 \(E_{ij}\) 读取向量的第 \(j\) 个坐标，把这个值写入第 \(i\) 个坐标，其余坐标置零。任何非零稳定子空间若含非零向量 \(v\)，取其一个非零坐标 \(v_j\)，则 \(E_{ij}v=v_j e_i\)；再乘 \(v_j^{-1}\)，便知它含有每个标准基向量，因而为全空间。这是不可约作用的一个基本例子。

% !TeX root = ../../Book2.tex
% 纲要标题：### 11.2 反环与包络代数
% 纲要位置：计划纲要.md，第 1414 行。

\section{反环与包络代数}
\label{b2:sec:1102}

左乘与右乘都可以表示为线性算子，但右乘算子的复合会颠倒乘法次序。反环记录这个次序，包络代数则把两个相容的作用合并成一个左作用。以下仍固定底域 \(k\)；所说的代数双模，其左右两侧的 \(k\) 标量作用须一致。

% 纲要要点（供目录核对）：
%
% * 反环与右作用的复合次序
% * 双模作为包络代数上的模，双边理想作为子模
% * 正则模自同态、矩阵包络作用及多项式商模型
%

\subsection{反环与左右模转换}
\label{b2:subsec:1102:01}

设 \(k\) 为域、\(A\) 为结合含幺 \(k\) 代数。取与 \(A\) 相同的向量空间和单位元，并规定乘法 \(a\cdot_{\mathrm{op}}b=ba\)，所得 \(k\) 代数称为 \(A\) 的反代数 \(A^{\mathrm{op}}\)。它就是\subsecref{b2:subsec:0201:01}所定义的反环，带上原有的 \(k\) 结构。将右 \(A\) 模 \(M\) 改写为左 \(A^{\mathrm{op}}\) 模时，规定 \(a^{\mathrm{op}}m=ma\)；原右作用的结合律保证这个左作用结合。

\ProseParagraphBreak
若 \(D_a(m)=ma\)，先右乘 \(b\)、再右乘 \(a\) 给出
\[
m\longmapsto mb\longmapsto(mb)a=m(ba),
\qquad D_a\circ D_b=D_{ba}.
\]
原来的右作用把乘法次序反转，改用反代数的乘法后，\(a^{\mathrm{op}}\mapsto D_a\) 就是通常的代数同态。例如 \(A=M_2(k)\) 时，先右乘 \(E_{12}\)、再右乘 \(E_{21}\)，得到 \(D_{E_{21}}\circ D_{E_{12}}=D_{E_{11}}\)；反向复合则为 \(D_{E_{12}}\circ D_{E_{21}}=D_{E_{22}}\)。

\subsection{双模作为包络代数上的模}
\label{b2:subsec:1102:02}

对两个 \(k\) 代数 \(A,B\)，\(A\otimes_kB\) 的乘法由
\[
(a\otimes b)(c\otimes d)=ac\otimes bd
\]
双线性延伸。四变量公式分别线性，故对两次张量关系良定义；在纯张量上，结合律由 \(A,B\) 的结合律得到，单位为 \(1_A\otimes1_B\)。这里两个代数本身都不必交换，只用到底域标量在各自中心。

\ProseParagraphBreak
对一个双模，我们希望用一个标量同时记录左乘 \(a\) 与右乘 \(b\)，使它作用为 \(m\mapsto amb\)。若先施加左乘 \(c\)、右乘 \(d\)，再施加左乘 \(a\)、右乘 \(b\)，连续作用为
\[
m\longmapsto cmd\longmapsto acmdb.
\]
左侧合成为 \(ac\)，右侧却合成为 \(db\)。因此记录右作用的因子须使用反代数，因为 \(b^{\mathrm{op}}d^{\mathrm{op}}=(db)^{\mathrm{op}}\)；张量积再把对这两个参数的双线性作用化为一个线性作用。

\begin{definition}\label{b2:def:enveloping-algebra}
设 \(k\) 为域、\(A\) 为结合含幺 \(k\) 代数。\(A\) 的\term[baoluo-daishu]{包络代数}[enveloping algebra]定义为
\[
A^e=A\otimes_k A^{\mathrm{op}}.
\]
若 \(k\) 向量空间 \(M\) 上有幺性的左、右 \(A\) 作用，两者对 \(k\) 双线性、满足 \((am)b=a(mb)\)，且两侧的底域标量作用都等于原有的 \(k\) 作用，则称 \(M\) 为 \(A\) 代数双模 \({}_AM_A\)。
\end{definition}
\symidx{\(A^e=A\otimes_k A^{\mathrm{op}}\)：代数的包络代数}

两侧标量一致的要求对应张量平衡关系 \((\lambda1_A)\otimes1=1\otimes(\lambda1_A)^{\mathrm{op}}\)。这两个张量是同一个元素，作用后分别得到 \((\lambda1_A)m\) 与 \(m(\lambda1_A)\)，所以它们必须相同，并等于原来的 \(\lambda m\)；仅有左右作用交换还不够。

\ProseParagraphBreak
这里的包络代数把代数的左、右作用合在一起；Lie 代数的泛包络代数 \(U(\mathfrak g)\) 则由 Lie 括号的另一种泛性质定义。

\begin{theorem}\label{b2:thm:enveloping-bimodules}
设 \(k\) 为域、\(A\) 为结合含幺 \(k\) 代数、\(A^e=A\otimes_kA^{\mathrm{op}}\)，\(M\) 为 \(k\) 向量空间。在 \(M\) 上指定与 \(k\) 标量作用相容的幺性 \(A\) 代数双模结构，等价于指定与该标量作用相容的幺性左 \(A^e\) 模结构。相应作用为
\begin{equation}\label{b2:eq:enveloping-action}
(a\otimes b^{\mathrm{op}})m=amb.
\end{equation}
这两种结构的同态也完全相同。
\end{theorem}
\begin{proof}
从双模出发，\((a,b^{\mathrm{op}})\mapsto(m\mapsto amb)\) 是到 \(\operatorname{End}_k(M)\) 的双线性映射，因此诱导 \(A^e\) 上的线性作用。连续作用两个纯张量给出
\[
(a\otimes b^{\mathrm{op}})\bigl((c\otimes d^{\mathrm{op}})m\bigr)
=acmdb.
\]
反代数中的乘法次序给出
\[
(a\otimes b^{\mathrm{op}})(c\otimes d^{\mathrm{op}})
=ac\otimes(db)^{\mathrm{op}},
\]
其作用也得到 \(acmdb\)。纯张量生成性保证乘法相容，\(1\otimes1\) 则作用为恒等。

反过来，令 \(am=(a\otimes1)m\)、\(mb=(1\otimes b^{\mathrm{op}})m\)。第一项保持乘法，第二项因反环而给出右结合律；两种纯张量 \(a\otimes1\) 与 \(1\otimes b^{\mathrm{op}}\) 交换，所以两侧作用相容。张量平衡关系给出
\[
(\lambda1_A)\otimes1=\lambda(1\otimes1)
=1\otimes(\lambda1_A)^{\mathrm{op}},
\]
而 \(1\otimes1\) 作用为恒等，故两侧的标量作用都等于 \(\lambda m\)。两次构造恢复原作用。同态同时与两侧作用相容，当且仅当与它们的复合及线性组合相容，亦即 \(A^e\) 线性。
\end{proof}

正则双模 \(A\) 本身由左右乘法给出。它的 \(A^e\) 子模恰是 \(A\) 的双边理想：对 \(a\otimes1\) 与 \(1\otimes b^{\mathrm{op}}\) 的作用稳定，分别就是对左乘与右乘稳定；反过来，这两种稳定性又保证对所有张量的作用稳定。因此研究 \(A\) 的双边理想可以转为研究正则 \(A^e\) 模 \(A\) 的子模。当 \(A\ne0\) 时，没有非零真双边理想，也就等价于这个 \(A^e\) 模没有非零真子模。

\ProseParagraphBreak
例如 \(A=k[t]\)，则 \(A^e\cong k[x,y]\)：\(x\) 记录左乘 \(t\)，\(y\) 记录右乘 \(t\)，包络代数把这两个作用分别记录。在正则双模上，这两个作用相同，因此 \(x-y\) 零化整个模。更一般的双模上，它们只须交换，不必相同。

\subsection{自同态代数与复合作用}
\label{b2:subsec:1102:03}

\begin{proposition}\label{b2:prop:regular-bimodule-endomorphism}
设 \(k\) 为域、\(A\) 为结合含幺 \(k\) 代数。对左正则模 \({}_AA\)，有 \(k\) 代数同构
\[
A^{\mathrm{op}}\longrightarrow\operatorname{End}_A({}_AA),\qquad
b^{\mathrm{op}}\longmapsto R_b,
\quad R_b(x)=xb.
\]
对正则双模，还有 \(\operatorname{End}_{A^e}(A)\cong Z(A)\)。
\end{proposition}
\begin{proof}
\cref{b2:prop:regular-hom-endomorphism}的公式给出每个左模自同态恰为 \(x\mapsto xb\)，其中 \(b=f(1)\) 唯一；这里它自动 \(k\) 线性。又 \(R_b\circ R_c=R_{cb}\)，正好给出反环同构。这样的映射同时保持右乘，当且仅当 \(ab=ba\) 对每个 \(a\in A\) 成立：把右线性等式 \(f(xa)=f(x)a\) 先取 \(x=1\) 得必要性，任意 \(x\) 的等式由同一交换性得到。因此双模自同态恰对应中心元素。中心交换，复合仍对应其通常乘法。
\end{proof}

对任意左 \(A\) 模 \(M\)，记 \(E=\operatorname{End}_A(M)\)。其中的映射与所有左标量作用交换。若要把这些自同态再作为 \(M\) 的右标量，系数代数须取 \(E^{\mathrm{op}}\)：定义 \(m\cdot f=f(m)\)，在反代数中有 \(f\cdot_{\mathrm{op}}g=g\circ f\)，因而
\[
m\cdot(f\cdot_{\mathrm{op}}g)
=g\bigl(f(m)\bigr)=(m\cdot f)\cdot g.
\]
这给出右作用的结合律，恒等映射给出单位作用，而 \(f(am)=a\cdot f(m)\) 保证它与左 \(A\) 作用相容。\chapref{b2:ch:14}的稠密性定理将采用这个右除环约定。

\begin{proposition}\label{b2:prop:regular-enveloping-models}
设 \(k\) 为域，\(n\ge1\)，对以下各 \(k\) 代数 \(C\) 记 \(C^e=C\otimes_k C^{\mathrm{op}}\)。对 \(A=M_n(k)\)，正则双模的作用给出含幺 \(k\) 代数同构
\[
\Phi:A^e\longrightarrow\operatorname{End}_k(A),
\qquad a\otimes b^{\mathrm{op}}\longmapsto(x\mapsto axb).
\]
对 \(B=k[t]\)，通过 \(x\mapsto t\otimes1\)、\(y\mapsto1\otimes t^{\mathrm{op}}\) 识别 \(B^e\cong k[x,y]\) 后，正则双模 \(B\) 作为左 \(B^e\) 模有同构
\[
k[x,y]/(x-y)\longrightarrow B,
\qquad F+(x-y)\longmapsto F(t,t).
\]
\end{proposition}
\begin{proof}
由\cref{b2:thm:enveloping-bimodules}应用于正则双模 \(A\)，\(\Phi\) 是含幺代数同态。设 \(E_{ij}\) 为 \(M_n(k)\) 的矩阵单位，对 \(1\le i,j,\ell,m,r,s\le n\)，有
\[
\Phi(E_{ij}\otimes E_{\ell m}^{\mathrm{op}})(E_{rs})
=E_{ij}E_{rs}E_{\ell m}
=\delta_{jr}\delta_{s\ell}E_{im}.
\]
因此这个算子把一个指定的基向量 \(E_{j\ell}\) 送到 \(E_{im}\)，把其余基向量送到零。让四个指标遍历，所得算子正是 \(\operatorname{End}_k(A)\) 的全部矩阵单位，组成一组基。来源中的 \(E_{ij}\otimes E_{\ell m}^{\mathrm{op}}\) 也是一组张量基，所以 \(\Phi\) 将一组基双射地送到另一组基，既单射又满射。

由于 \(B\) 交换，\(B^{\mathrm{op}}=B\)。张量基 \(t^i\otimes t^j\) 与单项式基 \(x^iy^j\) 一一对应，且乘法相符，所以得到所述 \(B^e\cong k[x,y]\)。在正则双模上，\(x,y\) 都作用为乘 \(t\)。于是取值映射
\[
\theta:k[x,y]\longrightarrow B,\qquad F\longmapsto F(t,t)
\]
是满射，也是左 \(k[x,y]\) 模同态，其中目标上的作用为 \(F\cdot p(t)=F(t,t)p(t)\)。显然 \(x-y\in\ker\theta\)。反向把 \(F(x,y)\) 看作 \(k[y][x]\) 中的多项式，除以首一多项式 \(x-y\)，可唯一写成
\[
F(x,y)=(x-y)Q(x,y)+R(y).
\]
若 \(\theta(F)=0\)，就有 \(R(t)=0\) 于多项式环 \(k[t]\)，因此 \(R\) 是零多项式，\(F\in(x-y)\)。这证明 \(\ker\theta=(x-y)\)，由\cref{b2:thm:module-first-isomorphism}得到所列模同构。
\end{proof}

% !TeX root = ../../Book2.tex
% 纲要标题：### 11.3 幂等元与分解
% 纲要位置：计划纲要.md，第 1420 行。

\section{幂等元与分解}
\label{b2:sec:1103}

一个投影满足连续作用两次与作用一次相同，即 \(p^2=p\)。在含幺代数 \(A\) 内，若 \(e^2=e\)，右乘映射 \(R_e(a)=ae\) 就满足 \(R_e^2=R_e\)，并将正则左模投到左理想 \(Ae\)。环中的幂等元由此产生模的分解，但未必分离整个代数的乘法；区别来自这个幂等元是否位于中心。

% 纲要要点（供目录核对）：
%
% * 正交幂等元的模分解与中心幂等元的代数分解
% * 角环、Peirce块及其双模结构
% * 右乘投影、角块与模同态复合
%

\subsection{正交幂等元及单位分解}
\label{b2:subsec:1103:01}

\begin{definition}\label{b2:def:orthogonal-idempotents}
设 \(k\) 为域、\(A\) 为含幺结合 \(k\) 代数。若 \(e\in A\) 满足 \(e^2=e\)，则称 \(e\) 为\term[mideng-yuan]{幂等元}[idempotent]。若两个幂等元 \(e,f\in A\) 满足 \(ef=fe=0\)，则称它们正交。若有限族 \(e_1,\ldots,e_r\in A\) 两两正交且 \(\sum_i e_i=1_A\)，则称 \(1_A=e_1+\cdots+e_r\) 为单位元的正交幂等元分解。
\end{definition}

\begin{proposition}\label{b2:prop:regular-orthogonal-decomposition}
设 \(k\) 为域、\(A\) 为含幺结合 \(k\) 代数，\(1_A=e_1+\cdots+e_r\) 是正交幂等元分解。它给出正则模的分解
\[
{}_AA=\bigoplus_i Ae_i,\qquad A_A=\bigoplus_i e_iA.
\]
若每个 \(e_i\) 还位于中心，则它们给出代数同构 \(A\cong\prod_i Ae_i\)，其中因子 \(Ae_i\) 的单位为 \(e_i\)。
\end{proposition}
\begin{proof}
对任意 \(a\)，有 \(a=\sum_i ae_i\)。若 \(\sum_i a_i e_i=0\)，右乘 \(e_j\) 只留下 \(a_j e_j\)，所以各项为零，左模分解成立。左乘的同样计算给出右模分解。

当各 \(e_i\) 在中心时，\(Ae_i\) 是双边理想和有单位的子代数，不同因子的乘积为零。映射 \(a\mapsto(ae_i)_i\) 保持加法、标量与乘法，并把单位送到各因子的单位族。其逆为求和映射，由刚证的直和分解互逆，得到代数同构。
\end{proof}

非中心幂等元只保证模分解。例如 \(A=M_2(k)=AE_{11}\oplus AE_{22}\) 是按列的左模分解，但 \(E_{11}\in AE_{11}\)、\(E_{12}\in AE_{22}\)，且 \(E_{11}E_{12}=E_{12}\ne0\)。这两个左模因子之间仍有非零乘积，不能成为直积代数中互相乘得零的两个因子。

\subsection{角环与矩阵分块}
\label{b2:subsec:1103:02}

设 \(k\) 为域、\(A\) 为含幺 \(k\) 代数、\(e\in A\) 为幂等元。集合 \(eAe\) 对加法、标量及乘法封闭，且以 \(e\) 为单位元；它称为 \(A\) 对应于 \(e\) 的\term[jiao-huan]{角环}[corner ring]，在当前情形也可称角代数。若 \(e\ne1_A\)，包含 \(eAe\subseteq A\) 不保持单位元，因而不是本书约定的含幺代数同态。

\ProseParagraphBreak
在 \(A=M_n(k)\) 中，若 \(e=\operatorname{diag}(I_s,0)\)、\(1\le s\le n\)，两侧乘 \(e\) 就只保留左上 \(s\times s\) 块，其余各块为零，所以 \(eAe\cong M_s(k)\)。对单位元的正交分解 \(1=\sum_i e_i\)，\(e_i a e_j\) 同样选出第 \((i,j)\) 块；下面的分解说明这种矩阵分块计算在一般代数中也成立。

\begin{proposition}[幂等元分块]\label{b2:prop:peirce-decomposition}
设 \(k\) 为域，\(A\) 为含幺结合 \(k\) 代数，\(1_A=e_1+\cdots+e_r\) 为正交幂等元分解，则
\begin{equation}\label{b2:eq:peirce-decomposition}
A=\bigoplus_{i,j}e_iAe_j
\end{equation}
作为向量空间成立。乘法满足 \((e_iAe_j)(e_sAe_t)=0\) 当 \(j\ne s\)；当 \(j=s\) 时，乘积属于 \(e_iAe_t\)。
每个块 \(e_iAe_j\) 通过 \(A\) 内的乘法成为 \((e_iAe_i,e_jAe_j)\) 双模，左右单位分别为 \(e_i\) 与 \(e_j\)。
\end{proposition}
\begin{proof}
把 \(a=1a1\) 展开，得到 \(a=\sum_{i,j}e_i a e_j\)。若这类分量之和为零，左乘 \(e_s\)、右乘 \(e_t\) 就只留下第 \((s,t)\) 个分量，所以分解唯一。两分量相乘时中间出现 \(e_je_s\)，它在不同指标时为零，相同指标时为 \(e_j\)，给出所述乘法规则。角环 \(e_iAe_i\) 左乘及 \(e_jAe_j\) 右乘都保持 \(e_iAe_j\)，并由结合律满足左右作用的相容性；对其中的 \(x\)，有 \(e_ix=x=xe_j\)，所以这些作用幺性。
\end{proof}

这种分块称为\term[Peirce-fenjie]{Peirce 分解}[Peirce decomposition]。\footnote{本杰明·皮尔斯（Benjamin Peirce，1809-1880），美国数学家、天文学家，在《Linear Associative Algebra》中以幂等元研究结合代数的结构。}因此可以把 \(a\) 记作分块矩阵 \((e_i a e_j)\)。乘积第 \((i,j)\) 块为 \(\sum_s(e_i a e_s)(e_s b e_j)\)，正是矩阵乘法式。不同块可能是不同的向量空间，不能未经论证就把它们全当作同一个系数环；非对角块作为双模连接两端的角环。

\ProseParagraphBreak
在 \(M_n(k)\) 中取 \(e_i=E_{ii}\)，各角环为 \(kE_{ii}\)，各块为 \(kE_{ij}\)，于是恢复通常的矩阵条目。若改取上三角代数 \(T_2(k)\)，右上块 \(e_1Ae_2=kE_{12}\) 非零，左下块 \(e_2Ae_1=0\)。两个角环都同构于 \(k\)，但整个代数不是 \(k\times k\)：非零右上块仍参与乘法。

\subsection{模分解与环中幂等元}
\label{b2:subsec:1103:03}

前面模论中已经见过投影的像与核分解；在\cref{b2:prop:finite-projective-summand}中，有限生成投射模还由有限秩自由模上的幂等自同态刻画。这里重新写出任意模的投影分解，是为了与环中幂等元产生的右乘投影对应起来。

\begin{proposition}\label{b2:prop:idempotent-module-decomposition}
对任意含幺环 \(R\) 和幺性左 \(R\) 模 \(M\)，满足 \(p^2=p\) 的 \(p\in\operatorname{End}_R(M)\) 给出
\[
M=\operatorname{im}p\oplus\ker p.
\]
反过来，每个指定的直和分解都给出投影幂等元。一组有限两两正交且和为恒等的投影，对应有限个子模的内直和分解。
\end{proposition}
\begin{proof}
写 \(m=p(m)+\bigl(m-p(m)\bigr)\)，第二项在核中；若 \(x=p(y)\) 又满足 \(p(x)=0\)，则 \(x=p^2(y)=p(x)=0\)，所以和为直和。反向取沿其余分量投到指定分量的映射，便满足平方为自身。有限族时，和为恒等保证生成，两两正交保证从一个分量到另一个的投影为零，从而分解唯一。
\end{proof}

在正则左模上，\cref{b2:prop:regular-bimodule-endomorphism}说明模自同态是右乘，所以幂等元 \(e\) 给出的投影为 \(a\mapsto ae\)，像为 \(Ae\)，核为 \(A(1-e)\)。左乘 \(L_e(a)=ea\) 虽也满足 \(L_e^2=L_e\)，却未必是左模自同态：左线性要求 \(e(ra)=r(ea)\)，取 \(a=1\) 就得到 \(er=re\) 对所有 \(r\in A\) 成立。反过来，\(e\) 在中心时此条件确实成立。取 \(A=M_2(k)\)、\(e=E_{11}\)、\(r=E_{21}\)，则
\[
L_e(r)=E_{11}E_{21}=0,
\qquad rL_e(1)=E_{21}E_{11}=E_{21}\ne0.
\]
因此左正则模中的投影须用右乘；不能仅凭算子平方等于自身，就把它当作左模投影。

\begin{proposition}\label{b2:prop:corner-endomorphism}
设 \(A\) 为含幺环，\(e,f\in A\) 为幂等元。取值映射给出交换群同构
\[
\operatorname{Hom}_A(Ae,Af)\longrightarrow eAf,
\qquad h\longmapsto h(e).
\]
它的逆把 \(b\in eAf\) 送到 \(h_b:Ae\rightarrow Af\)、\(h_b(ae)=ab\)。若 \(g\in A\) 也是幂等元，\(b\in eAf\)、\(c\in fAg\)，则
\[
h_c\circ h_b=h_{bc}:Ae\longrightarrow Ag.
\]
特别地，有含幺环同构
\[
(eAe)^{\mathrm{op}}\longrightarrow\operatorname{End}_A(Ae),
\qquad b^{\mathrm{op}}\longmapsto h_b.
\]
若 \(k\) 为域、\(A\) 为含幺结合 \(k\) 代数，则上述取值同构为 \(k\) 线性同构，上述自同态环同构为含幺 \(k\) 代数同构。
\end{proposition}
\begin{proof}
若 \(h:Ae\to Af\) 左线性，因为 \(h(e)\in Af\)，可写成 \(h(e)=a_0f\)。又 \(h(e)=h(e^2)=e h(e)\)，所以
\[
h(e)=ea_0f\in eAf.
\]
对任意 \(a\in A\)，左线性又迫使 \(h(ae)=a\cdot h(e)\)，故 \(e\) 的像决定整个同态。

反过来，取 \(b\in eAf\)。若 \(ae=a'e\)，则 \((a-a')e=0\)，而 \(eb=b\)，所以 \((a-a')b=(a-a')eb=0\)，证明 \(h_b(ae)=ab\) 与代表元选择无关。由于 \(bf=b\)，有 \(ab\in Af\)；结合律给出左 \(A\) 线性。取值映射与 \(b\mapsto h_b\) 都保持加法，而 \(h_b(e)=eb=b\)，所以这两个加法群映射互为逆。

对 \(b\in eAf\)、\(c\in fAg\)，有 \(bc\in eAg\)，并且
\[
(h_c\circ h_b)(ae)=h_c(ab)=abc=h_{bc}(ae).
\]
这里 \(ab=(ab)f\) 作为 \(Af\) 中的元素，正可代入 \(h_c\) 的定义。取 \(f=g=e\)，便有 \(h_b\circ h_c=h_{cb}\)，所以参数乘法相对于通常的自同态复合恰好反转，得到 \((eAe)^{\mathrm{op}}\) 到自同态环的同构；\(h_e(ae)=ae\)，故角环的单位 \(e\) 对应恒等映射。当 \(A\) 还是 \(k\) 代数时，取值与其逆都保持 \(k\) 标量，故这些对应分别是所述的 \(k\) 线性同构与 \(k\) 代数同构。
\end{proof}

这个取值对应说明，Peirce 块 \(eAf\) 同时记录从左理想 \(Ae\) 到 \(Af\) 的所有左模同态，块乘法 \((eAf)(fAg)\subseteq eAg\) 对应先施加 \(h_b\)、再施加 \(h_c\) 的复合。取 \(f=e\)，便把正则模的分解细化到某个直和因子内部。后续判别 \(Ae\) 是否还能分解时，考察的是它的自同态环中有无非平凡幂等元，也就是角环 \(eAe\) 中的同一问题。

% !TeX root = ../../Book2.tex
% 纲要标题：### 11.4 中心化子
% 纲要位置：计划纲要.md，第 1426 行。

\section{中心化子}
\label{b2:sec:1104}

一个作用确定以后，可以再问哪些线性算子与它交换。这些算子组成中心化子；再取一次中心化子，则至少能找回原来的作用。但是否恰好找回原代数，是一个需要证明的问题。本节先给出完全可算的例子，再说明它与张量表示的联系。

% 纲要要点（供目录核对）：
%
% * 模自同态与中心化子，双中心化子的包含与幂等性
% * 正则左右作用及有限维张量因子的中心化子
% * 三角反例与 Schur–Weyl 对偶的动机及证明归属
%
% ---
%

\subsection{交换作用与中心化子}
\label{b2:subsec:1104:01}

设 \(k\) 为域、\(V\) 为 \(k\) 向量空间，\(\mathcal A\subseteq\operatorname{End}_k(V)\) 为含恒等算子的子代数。沿用\subsecref{b2:subsec:0804:03}的中心化子定义，记
\[
\mathcal A'=\bigl\{\,T\in\operatorname{End}_k(V):Ta=aT\text{ 对所有 }a\in\mathcal A\,\bigr\}.
\]
再记 \(\mathcal A''=(\mathcal A')'\)，称为\term[shuang-zhongxinhua-zi]{双中心化子}[double centralizer]。这些撇号在本节表示中心化子，不表示导数或对偶。
\symidx{\(\mathcal A'\)、\(\mathcal A''\)：给定作用的中心化子及双中心化子}

\begin{proposition}\label{b2:prop:commutant-endomorphism}
设 \(k\) 为域、\(A\) 为结合含幺 \(k\) 代数、\(V\) 为 \(k\) 向量空间，\(\rho:A\to\operatorname{End}_k(V)\) 为含幺表示。对 \(\mathcal C\subseteq\operatorname{End}_k(V)\)，用 \(\mathcal C'\) 表示 \(\operatorname{End}_k(V)\) 中与 \(\mathcal C\) 全部元素交换的算子所成的中心化子，则
\[
\rho(A)'=\operatorname{End}_A(V),\qquad \rho(A)\subseteq\rho(A)''.
\]
若 \(\mathcal A\subseteq\mathcal B\subseteq\operatorname{End}_k(V)\)，则
\[
\mathcal B'\subseteq\mathcal A',\qquad
\mathcal A''\subseteq\mathcal B''.
\]
对任意算子族 \(\mathcal C\subseteq\operatorname{End}_k(V)\)，还有
\[
\mathcal C\subseteq\mathcal C'',\qquad
\mathcal C'''=\mathcal C',\qquad
\mathcal C''''=\mathcal C''.
\]
\end{proposition}
\begin{proof}
\cref{b2:prop:algebra-modules-representations}中模同态的相容公式在定义域、陪域均为 \(V\) 时，就是 \(T\rho(a)=\rho(a)T\)。在向量 \(v\) 上，这个等式写成 \(T(av)=a\cdot T(v)\)，故第一个等式成立。任意 \(c\in\mathcal C\) 都与 \(\mathcal C'\) 的全部元素交换，所以 \(\mathcal C\subseteq\mathcal C''\)，这也给出 \(\rho(A)\subseteq\rho(A)''\)。

若一个算子与 \(\mathcal B\) 的全部元素交换，它当然与子集 \(\mathcal A\) 的全部元素交换，故 \(\mathcal B'\subseteq\mathcal A'\)。再取一次中心化子并反转包含，得到 \(\mathcal A''\subseteq\mathcal B''\)。由 \(\mathcal C\subseteq\mathcal C''\) 得 \(\mathcal C'''\subseteq\mathcal C'\)；把扩大性质用于 \(\mathcal C'\)，又得 \(\mathcal C'\subseteq\mathcal C'''\)。于是 \(\mathcal C'''=\mathcal C'\)，两侧再取中心化子即得 \(\mathcal C''''=\mathcal C''\)。
\end{proof}

与更多算子交换会施加更多限制，所以中心化子反转包含；连续取两次后又保持包含。上述等式还说明，双中心化子扩大原算子族，并且再取双中心化子不会继续增大。这是它作为一种闭包运算的意义。

\ProseParagraphBreak
若两个代数在同一空间上的作用彼此交换，其中一个便落在另一个的中心化子中。只有证明反向包含，才能称它们互为中心化子。维数有限时可以把交换条件写成线性方程组；这通常比直接猜测所有相容映射更可靠。

\subsection{双中心化子的基本例子}
\label{b2:subsec:1104:02}

\begin{theorem}[正则作用的双中心化子]\label{b2:thm:regular-double-centralizer}
设 \(k\) 为域、\(A\) 为结合含幺 \(k\) 代数。在 \(\operatorname{End}_k(A)\) 中，令
\[
L(A)=\{L_a:x\mapsto ax\mid a\in A\},\qquad
R(A)=\{R_a:x\mapsto xa\mid a\in A\}.
\]
映射 \(a\mapsto L_a\) 与 \(a^{\mathrm{op}}\mapsto R_a\) 分别给出 \(k\) 代数同构 \(A\cong L(A)\) 与 \(A^{\mathrm{op}}\cong R(A)\)。两个作用代数互为中心化子，因此 \(L(A)''=L(A)\)，不要求 \(A\) 有限维。
\end{theorem}
\begin{proof}
复合满足 \(L_aL_b=L_{ab}\)、\(R_aR_b=R_{ba}\)，所以右乘的复合顺序对应反代数乘法。两种映射都线性、保单位元，并且算子在 \(1\) 处的值恢复 \(a\)，故它们单射；由像的定义，它们满射到相应作用代数。

\(T\) 与全部左乘交换，当且仅当它左 \(A\) 线性。正则左模由 \(1\) 生成；\cref{b2:prop:regular-bimodule-endomorphism}给出这样的 \(T\) 由 \(b=T(1)\) 唯一确定，且 \(T(x)=xb\)，故 \(L(A)'=R(A)\)。若 \(T\) 与全部右乘交换，则对任意 \(a\)，
\[
T(a)=T\bigl(R_a(1)\bigr)=R_a\bigl(T(1)\bigr)=T(1)a.
\]
整个算子仍由 \(T(1)\) 决定，这次它是左乘 \(T(1)\)。反向左乘与右乘由结合律交换，得到 \(R(A)'=L(A)\)。
\end{proof}

另一个重要例子来自两个向量空间的张量。令 \(U,W\) 为非零有限维空间。代数 \(\operatorname{End}_k(U)\) 作用于第一因子，\(\operatorname{End}_k(W)\) 作用于第二因子，两种作用交换。选定 \(U\) 的一组 \(m\) 元基后，\(U\otimes_kW\) 是 \(W\) 的 \(m\) 个副本的直和；第一因子的矩阵作用调配这些副本。与所有这些调配交换的算子，必须在每个副本上作同一个变换。

\begin{proposition}\label{b2:prop:tensor-factor-centralizers}
设 \(k\) 为域，\(U,W\) 为非零有限维 \(k\) 向量空间。在 \(\operatorname{End}_k(U\otimes_kW)\) 中，
\[
\operatorname{End}_k(U)\otimes k\operatorname{id}_W,
\qquad k\operatorname{id}_U\otimes\operatorname{End}_k(W)
\]
互为中心化子。此外，映射
\[
\begin{aligned}
&\operatorname{End}_k(U)\otimes_k\operatorname{End}_k(W)
\longrightarrow\operatorname{End}_k(U\otimes_kW),\\
&F\otimes G\longmapsto
\bigl(u\otimes w\mapsto F(u)\otimes G(w)\bigr)
\end{aligned}
\]
是 \(k\) 代数同构。
\end{proposition}
\begin{proof}
选 \(U\) 的基 \(u_1,\ldots,u_m\)。任意线性算子 \(T\) 可唯一写成
\[
T(u_j\otimes w)=\sum_i u_i\otimes T_{ij}(w),
\qquad T_{ij}\in\operatorname{End}_k(W).
\]
\(E_{aa}\otimes\operatorname{id}_W\) 是到第 \(a\) 个副本 \(ku_a\otimes W\) 的投影。若 \(T\) 与这些投影交换，它就保留各个副本，故 \(T_{ij}=0\) 当 \(i\ne j\)。而 \(E_{ab}\otimes\operatorname{id}_W\) 把第 \(b\) 个副本移到第 \(a\) 个副本，在其中保留 \(w\) 不变。与这个算子交换，便要求两个副本上的变换相同：在 \(u_b\otimes w\) 上比较，得到 \(T_{aa}(w)=T_{bb}(w)\)。因此 \(T=\operatorname{id}_U\otimes S\)，其中 \(S\) 为公共对角块。反向这样的算子确实交换，得到第一个中心化子。交换两因子的角色，得到另一个。

再选 \(W\) 的基 \(w_1,\ldots,w_\ell\)。张量 \(E_{ij}\otimes E_{ab}\) 将基向量 \(u_j\otimes w_b\) 送到 \(u_i\otimes w_a\)，将其余基向量送到零，所以这些张量恰对应全部矩阵单位。它们在两边均为基，故作用映射是线性同构；逐因子复合保证它保持乘法与单位。
\end{proof}

例如 \(k^n\oplus k^n\cong k^n\otimes_k k^2\)（\(n\ge1\)）时，\(\operatorname{diag}(B,B)\)、\(B\in M_n(k)\)，作用于第一因子。其中心化子为 \(\operatorname{id}_{k^n}\otimes\operatorname{End}_k(k^2)\)，允许在两个相同的 \(k^n\) 副本之间任意混合，而在各副本内部保留原来的矩阵作用。\cref{b2:ex:11-repeated-matrix-blocks}将这个结论写成具体的分块矩阵。

\ProseParagraphBreak
这里自同态代数的张量同构不能无条件推广到无限维空间。取 \(U,W\) 分别以 \(u_1,u_2,\ldots\) 与 \(w_1,w_2,\ldots\) 为基，定义
\[
Q(u_i\otimes w_j)=
\begin{cases}
u_i\otimes w_i,&i=j,\\
0,&i\ne j.
\end{cases}
\]
这在张量积的一组基上规定了像，故给出线性算子。若 \(Q=\sum_{r=1}^s A_r\otimes B_r\)，记 \(\alpha_i:U\to k\) 为第 \(i\) 个坐标泛函，对输入 \(u_i\otimes w\) 的输出取第 \(i\) 个 \(U\) 坐标，就得到
\[
P_i(w)=\sum_{r=1}^s\alpha_i(A_ru_i)B_r(w),
\]
其中 \(P_i\in\operatorname{End}_k(W)\) 保留 \(w_i\) 坐标，并把其余坐标归零。因此所有 \(P_i\) 都在有限维空间 \(\operatorname{span}_k\{B_1,\ldots,B_s\}\) 中。但这些投影线性无关：任意有限线性关系分别作用于对应的 \(w_i\)，就迫使每个系数为零。矛盾。所以这个 \(Q\) 不在张量作用映射的像中。

\ProseParagraphBreak
双中心化子等式不能推广到每个子代数。取 \(\mathcal A=T_2(k)\subseteq M_2(k)\)。与 \(E_{11}\) 交换迫使矩阵为对角矩阵，再与 \(E_{12}\) 交换迫使两个对角元相等。因此 \(\mathcal A'=kI_2\)，而 \(\mathcal A''=M_2(k)\supsetneq\mathcal A\)。原代数保留非零真子空间 \(ke_1\)，故作用不是不可约的。这就给出了从中心化子仅含标量反推不可约性的反例。

\subsection{Schur–Weyl 对偶的环论背景}
\label{b2:subsec:1104:03}

\chapref{b2:ch:08}已在 \(V^{\otimes d}\) 上构造一般线性群的对角作用与 \(S_d\) 的置换作用，并在\cref{b2:prop:commuting-tensor-actions}中证明它们交换。因而两种作用的线性生成代数互相包含于对方的中心化子。对于有限维复空间，\cref{b2:claim:schur-weyl-preview}给出它们互为中心化子的等式；证明见第五卷第7.3节。

\ProseParagraphBreak
这里的情形与\cref{b2:prop:tensor-factor-centralizers}不同。后者允许独立选取 \(F\in\operatorname{End}_k(U)\) 与 \(G\in\operatorname{End}_k(W)\)，作用为 \(u\otimes w\mapsto F(u)\otimes G(w)\)；Schur--Weyl 情形则让同一个 \(g\) 同时作用在所有因子上，得到 \(v_1\otimes\cdots\otimes v_d\mapsto gv_1\otimes\cdots\otimes gv_d\)，另一个作用改变因子的排列。独立因子上的中心化子计算没有计算出同一个 \(g\) 的对角作用的全部中心化子，因而不能代替 Schur--Weyl 对偶中反向包含的证明。

\ProseParagraphBreak
若只考察交换两个因子的算子 \(\tau\)，在特征不为二的域上，第八章的两个投影 \((\operatorname{id}\pm\tau)/2\) 把张量积分成对称与交替两部分。与 \(\tau\) 交换的算子恰好分别保持这两部分，因而可以在每一部分内部独立作用。\cref{b2:ex:11-tensor-flip-centralizer}在 \(V=k^2\) 上计算这一中心化子及其双中心化子。

\ProseParagraphBreak
\chapref{b2:ch:14}将研究不可约作用的中心化子及稠密性；有限维且底域代数闭时，不可约性会迫使作用代数等于整个自同态代数。上三角例子则表明，只知中心化子由标量组成还不足以作这个推断。双中心化子问题的答案取决于整个作用，而不只取决于原代数的抽象乘法表。

\ProseParagraphBreak
若另具备Hilbert空间和有界算子的基础，可参阅\appref{b2:app:D}的\cref{b2:thm:operator-bicommutant}。该结果在含幺且对伴随封闭的条件下，把双交换子与强、弱算子闭包联系起来；这里的纯代数中心化子计算并不自动包含那些拓扑条件。


\begin{chaptersummary}
有限维代数在选定一组有限基后，其结构常数记录全部乘法，结合律与单位条件成为有限组方程；换基会改变这些常数，却保留代数结构。把代数作用写成到自同态代数的同态，可以同时处理群表示、矩阵作用和一般模。

\ProseParagraphBreak
右作用对应反环，两个相容的左右作用对应包络代数的一个左作用，双边理想因而成为正则包络代数模的子模。正则左模的自同态来自右乘，正则双模的自同态则由中心决定。环中的正交幂等元给出正则模与角块分解；角块 \(eAf\) 记录 \(Ae\to Af\) 的模同态，块乘法记录这些同态的复合。只有幂等元还满足中心性时，这种模分解才进一步给出代数乘积。

\ProseParagraphBreak
中心化子描述与一个作用交换的全部算子，再取中心化子总包含原作用代数，却未必等于它。正则作用及有限维张量因子的独立作用具有双中心化子等式，上三角作用则说明标量中心化子本身不足以保证不可约。后续结构理论要补充的，正是这些具体计算背后的有限性与模分解条件。
\end{chaptersummary}

\BookExercises

\begin{exercise}\label{b2:ex:11-two-dimensional-constants}
设 \(k\) 为域，非零结合含幺 \(k\) 代数 \(A\) 有基 \(1,x\)，其中 \(x^2=\alpha x+\beta\)、\(\alpha,\beta\in k\)。令 \(y=ux+v\)，其中 \(u\in k^\times\)、\(v\in k\)。求基 \(1,y\) 下的乘法关系，并证明 \(\Delta=\alpha^2+4\beta\) 在这种换基下变为 \(u^2\Delta\)。说明等式本身是否要求特征不为二。
\end{exercise}
\begin{solution}
展开并消去 \(x=u^{-1}(y-v)\)，得到
\[
y^2=(u\alpha+2v)y+(u^2\beta-u\alpha v-v^2).
\]
新系数为 \(\alpha'=u\alpha+2v\)、\(\beta'=u^2\beta-u\alpha v-v^2\)，所以
\[
(\alpha')^2+4\beta'=u^2\alpha^2+4u^2\beta=u^2\Delta.
\]
全部计算都是整系数恒等式，不要求特征不为二。这个变换律说明换基只会把 \(\Delta\) 乘以一个非零平方，但分类还须证明相应的标准形：特征不为二时，取 \(z=x-\alpha/2\) 得 \(z^2=\Delta/4\)，再分别辨认 \(\Delta=0\)、非零平方及非平方三种情形；特征二时，\(\Delta=\alpha^2\) 甚至不记录 \(\beta\)。另有 \(A\cong k[t]/(t^2-\alpha t-\beta)\)：送 \(t\) 到 \(x\) 给出满同态，商中 \(1,t\) 与目标的 \(1,x\) 分别为基，所以它单射。
\end{solution}

\begin{exercise}\label{b2:ex:11-associative-table}
设 \(k\) 为域、\(\lambda,\mu\in k\)。在有基 \(1,x,y\) 的 \(k\) 向量空间中，令 \(1\) 为单位，并规定交换的基乘法
\(x^2=y\)、\(xy=yx=\lambda y\)、\(y^2=\mu y\)。求这一双线性乘法满足结合律的充要条件，并在满足条件时给出单生成元商代数描述。
\end{exercise}
\begin{solution}
由 \((x^2)y=x(xy)\)，必须有 \(\mu y=\lambda^2y\)，所以 \(\mu=\lambda^2\)。若该条件成立，候选的单生成元关系是 \(t^3=\lambda t^2\)。考察结合代数 \(B=k[t]/(t^3-\lambda t^2)\)，它以 \(1,t,t^2\) 为基，且 \(t\cdot t^2=\lambda t^2\)、\((t^2)^2=\lambda^2t^2\)。基对应 \(1\mapsto1,t\mapsto x,t^2\mapsto y\) 因而保留所有给定基乘法，使 \(B\) 与题给双线性乘法结构同构。后者由此结合，且得到 \(A\cong B\)。这种构造也核对了充分性，而不只检查某一个必要等式。
\end{solution}

\begin{exercise}\label{b2:ex:11-nilpotent-representation}
设 \(k\) 为域、\(A=k[t]/(t^3)\)。在 \(V=k^3\) 上分别令 \(t\) 的剩余类作用为三阶幂零 Jordan 块 \(J_3\)，以及 \(J_2\oplus0\)。证明这两者都给出 \(A\) 的表示，并求表示的核。
\end{exercise}
\begin{solution}
两个算子都满足 \(T^3=0\)，所以多项式代入 \(f\mapsto f(T)\) 经过商代数 \(A\)，且保持乘法与单位。对 \(J_3\)，若 \(aI+bJ_3+cJ_3^2=0\)，依次比较主对角线、第一条及第二条上对角线，得到 \(a=b=c=0\)，所以表示忠实。对 \(J_2\oplus0\)，有 \(T\ne0,T^2=0\)，且 \(I,T\) 线性无关，所以代入核正好为 \(k\bar t^2=(\bar t^2)\)。同一代数作用在相同维数空间上，仍可有不同的零化子。
\end{solution}

\begin{exercise}\label{b2:ex:11-group-algebra-universal}
设 \(k\) 为域，\(G\) 为群，\(N\triangleleft G\)。令 \(I\) 为 \(k[G]\) 中由 \(n-1\)、\(n\in N\) 生成的双边理想。证明商群映射诱导代数同构
\[
k[G]/I\cong k[G/N].
\]
说明在群代数中加入这些关系，为什么正好把同一 \(N\) 陪集中的基元素识别，而不会产生额外的线性关系。
\end{exercise}
\begin{solution}
由\cref{b2:prop:group-algebra-representations}，\(g\mapsto gN\) 唯一延拓为满代数同态 \(q:k[G]\to k[G/N]\)。每个 \(n-1\) 都在核中，故 \(I\subseteq\ker q\)。若有限和 \(z=\sum_g a_gg\) 在核中，由于各陪集是 \(k[G/N]\) 的基，每个陪集内的系数和都为零。对与 \(z\) 的支集相交的每个陪集 \(C\)，取代表元 \(r_C\)，便有
\[
\sum_{g\in C}a_gg
=\sum_{g\in C}a_g(g-r_C).
\]
这里所有和都只取原有限支集中的元素。对 \(g\in C=r_CN\)，写 \(g=r_Cn\)，则 \(g-r_C=r_C(n-1)\in I\)。逐陪集相加得 \(z\in I\)，所以 \(\ker q=I\)，商代数同构随之成立。陪集内系数和为零恰好描述了这些识别关系的全部线性后果。
\end{solution}

\begin{exercise}\label{b2:ex:11-augmentation-generators}
设 \(k\) 为域，群 \(G\) 由子集 \(X\) 生成，\(\epsilon:k[G]\to k\) 为增广同态。证明增广理想 \(\ker\epsilon\) 作为左理想已由 \(x-1\)、\(x\in X\) 生成，不需要把所有群元素逐个列入生成元。说明同样结论也适用于右理想生成。
\end{exercise}
\begin{solution}
令 \(L\) 为这些元素生成的左理想。恒等式
\[
gh-1=(g-1)+g(h-1),\qquad x^{-1}-1=-x^{-1}(x-1)
\]
说明若已知两个群元素减一属于 \(L\)，则其乘积减一也属于 \(L\)，并且每个生成元的逆减一也在 \(L\)。对生成字长度归纳，得到所有 \(g-1\in L\)。若 \(\sum_g a_gg\in\ker\epsilon\)，其系数和为零，故它等于 \(\sum_g a_g(g-1)\)，也属于 \(L\)。这证明 \(\ker\epsilon\subseteq L\)；反向包含由 \(\epsilon(x-1)=0\) 得到。右理想版本使用 \(gh-1=(g-1)h+(h-1)\) 与 \(x^{-1}-1=-(x-1)x^{-1}\)，其余归纳相同。
\end{solution}

\begin{exercise}\label{b2:ex:11-cyclic-two-units}
设 \(k\) 为域、\(C_2=\{1,g\}\)、\(g^2=1\)。对 \(a,b\in k\)，求 \(k[C_2]\) 中 \(a+bg\) 可逆的充要条件和逆元公式，允许 \(\operatorname{char}k=2\)。
\end{exercise}
\begin{solution}
左乘 \(a+bg\) 在基 \(1,g\) 下的矩阵为 \(\begin{pmatrix}a&b\\b&a\end{pmatrix}\)。若元素可逆，该矩阵可逆，故 \(a^2-b^2\ne0\)。反过来，若此数非零，直接相乘得到
\[
(a+bg)^{-1}=\frac{a-bg}{a^2-b^2}.
\]
这是双侧逆，所以条件充分。在特征二中，条件变为 \((a+b)^2\ne0\)，即 \(a+b\ne0\)，公式仍成立。可逆元不限于非零标量乘群元素，例如在 \(\mathbb Q[C_2]\) 中 \(2+g\) 也可逆。
\end{solution}

\begin{exercise}\label{b2:ex:11-matrix-opposite}
设 \(k\) 为域、\(n\ge1\)。证明转置给出代数同构 \(M_n(k)^{\mathrm{op}}\to M_n(k)\)，但在 \(n\ge2\) 时不是从 \(M_n(k)\) 到自身的代数同态。用两个矩阵单位写出具体失败等式。
\end{exercise}
\begin{solution}
来源反环中 \(A\cdot_{\mathrm{op}}B=BA\)，所以
\((A\cdot_{\mathrm{op}}B)^{\mathsf t}=A^{\mathsf t}B^{\mathsf t}\)。转置线性、保持单位且自身为逆，故是所述同构。在原乘法下取 \(A=E_{12},B=E_{21}\)，则 \((AB)^{\mathsf t}=E_{11}\)，而 \(A^{\mathsf t}B^{\mathsf t}=E_{21}E_{12}=E_{22}\)，两者不同。
\end{solution}

\begin{exercise}\label{b2:ex:11-matrix-enveloping}
设 \(k\) 为域，\(A=k\times k\)，\(e_1=(1,0)\)、\(e_2=(0,1)\)。计算正则双模作用 \(\Phi:A^e\to\operatorname{End}_k(A)\) 的像与核，其中 \(\Phi(a\otimes b)(x)=axb\)。说明左正则作用与右正则作用各自忠实，为什么仍不足以使整个包络代数的作用忠实。
\end{exercise}
\begin{solution}
因为 \(A\) 交换，\(A^{\mathrm{op}}=A\)，而 \(e_i\otimes e_j\)、\(1\le i,j\le2\) 构成 \(A^e\) 的基。由 \(e_ie_j=\delta_{ij}e_i\)，有
\[
\Phi(e_i\otimes e_j)(e_\ell)
=\delta_{i\ell}\delta_{j\ell}e_\ell.
\]
所以 \(e_1\otimes e_1\)、\(e_2\otimes e_2\) 给出两个坐标投影，非对角的两个张量作用为零。由此
\[
\operatorname{im}\Phi
=\left\{\begin{pmatrix}a&0\\0&b\end{pmatrix}:a,b\in k\right\},
\qquad
\ker\Phi=k(e_1\otimes e_2)\oplus k(e_2\otimes e_1).
\]
左乘或右乘 \(a\) 若为零，在 \(1\) 上求值便得 \(a=0\)，故两种单侧作用各自忠实。但包络代数把两侧分别记录，还含有上述非零张量，单侧的忠实性没有排除它们。与\cref{b2:prop:regular-enveloping-models}的矩阵代数相比，这里的左右乘有限和只能产生对角算子。
\end{solution}

\begin{exercise}\label{b2:ex:11-one-dimensional-bimodules}
设 \(k\) 为域，\(A=k[t]\)，\(\alpha,\beta,\gamma,\delta\in k\)。在一维空间 \(M_{\alpha,\beta}=k\) 上规定左乘 \(t\) 为乘 \(\alpha\)，右乘 \(t\) 为乘 \(\beta\)，两侧 \(k\) 标量都按通常方式作用。证明它总是 \(A\) 双模，并求
\(\operatorname{Hom}_{A^e}(M_{\alpha,\beta},M_{\gamma,\delta})\)。
\end{exercise}
\begin{solution}
左右作用分别为 \(f(t)m=f(\alpha)m\)、\(m\cdot g(t)=g(\beta)m\)。它们都保持多项式乘法、单位与底域标量，且两个标量作用交换，所以确为双模。任一线性映射是乘 \(c\in k\)，保持左、右作用分别要求 \(c\alpha=\gamma c\)、\(c\beta=\delta c\)。若两对参数相同，任意 \(c\) 都可行，Hom 为 \(k\)；只要有一项不同，域中消去非零差便得 \(c=0\)，Hom 为零。两个作用可以不同，这不违反双模的相容性。
\end{solution}

\begin{exercise}\label{b2:ex:11-dual-number-enveloping-matrices}
设 \(k\) 为域，\(A=k[\varepsilon]/(\varepsilon^2)\)，\(M=k^3\)，并以 \(E_{ij}\) 表示 \(M_3(k)\) 的矩阵单位。规定 \(\varepsilon\) 在 \(M\) 上的左作用为 \(E_{12}\)、右作用为 \(E_{13}\)，两侧标量均按通常方式作用。证明这些规定使 \(M\) 成为 \(A\) 双模。将 \(A^e\) 识别为 \(k[x,y]/(x^2,y^2)\)，写出 \(1,x,y,xy\) 在 \(M\) 上的表示矩阵，并说明如何从这些矩阵恢复左右 \(A\) 作用。
\end{exercise}
\begin{solution}
矩阵 \(E_{12}\)、\(E_{13}\) 的平方均为零，且 \(E_{12}E_{13}=E_{13}E_{12}=0\)，故它们分别满足 \(\varepsilon^2=0\) 并彼此交换。这给出相容的左、右幺性作用。令 \(x=\varepsilon\otimes1\)、\(y=1\otimes\varepsilon^{\mathrm{op}}\)。因 \(A\) 交换，\(A^e\) 由交换的 \(x,y\) 生成，满足 \(x^2=y^2=0\)；两边都以 \(1,x,y,xy\) 为 \(k\) 基，所以 \(A^e\cong k[x,y]/(x^2,y^2)\)。这四个元素在 \(M\) 上依次表示为 \(I_3,E_{12},E_{13},0\)。把所得表示分别限制到 \(A\otimes1\) 与 \(1\otimes A^{\mathrm{op}}\)，便恢复左、右 \(A\) 作用。
\end{solution}

\begin{exercise}\label{b2:ex:11-polynomial-bimodule-annihilator}
设 \(k\) 为域，\(n\ge2\)，\(A=k[t]\)，\(M=k[t]/(t^n)\) 以商中的乘法作左右 \(A\) 作用。把 \(A^e\) 识别为 \(k[x,y]\)，其中 \(x,y\) 分别记录左右乘 \(t\)。证明
\[
\operatorname{Ann}_{A^e}(M)=(x-y,y^n),
\qquad M\cong k[x,y]/(x-y,y^n)
\]
作为 \(A^e\) 模成立。说明为何也可把 \(y^n\) 换成 \(x^n\)，并求 \(\dim_kM\)。
\end{exercise}
\begin{solution}
用 \(x-y\) 对 \(F\in k[y][x]\) 作首一带余除法，写成 \(F=(x-y)Q+r(y)\)。在 \(M\) 上左右乘 \(t\) 相同，故 \(F\) 作用为乘以 \(r(t)+(t^n)\)。它零化 \(M\) 当且仅当零化 \(\bar1\)，即 \(t^n\mid r(t)\)。因此零化子正好是 \((x-y,y^n)\)。求值映射 \(F\mapsto F(t,t)+(t^n)\) 满射到 \(M\)，并保持 \(A^e\) 作用，给出所求商模同构。又
\[
x^n-y^n=(x-y)(x^{n-1}+x^{n-2}y+\cdots+y^{n-1}),
\]
所以 \((x-y,y^n)=(x-y,x^n)\)。商中 \(\bar1,\bar t,\ldots,\bar t^{n-1}\) 为基，故维数为 \(n\)。与正文正则双模的关系 \(x=y\) 相比，截断还加入了 \(t^n=0\)。
\end{solution}

\begin{exercise}\label{b2:ex:11-triangular-idempotents}
设 \(k\) 为域。求上三角代数 \(T_2(k)\) 的全部幂等元，并从中辨认中心幂等元。由此判断它能否分解为两个非零代数的直积。
\end{exercise}
\begin{solution}
写 \(e=\begin{pmatrix}a&b\\0&c\end{pmatrix}\)。\(e^2=e\) 等价于 \(a^2=a,c^2=c,(a+c-1)b=0\)。域中 \(a,c\in\{0,1\}\)。若 \(a=c\)，则 \(b=0\)，得到零与单位；若 \((a,c)=(1,0)\) 或 \((0,1)\)，则 \(b\) 任意。因此其余幂等元为
\[
\begin{pmatrix}1&b\\0&0\end{pmatrix},\qquad
\begin{pmatrix}0&b\\0&1\end{pmatrix}\quad(b\in k).
\]
与 \(E_{11}\) 交换要求 \(b=0\)，再与 \(E_{12}\) 交换要求 \(a=c\)，所以中心幂等元只有零与单位。若存在非零乘积分解，第一因子的单位 \((1,0)\) 会给出非平凡中心幂等元，矛盾。因此不存在这种分解，虽然正则模仍可按 \(E_{11},E_{22}\) 分成两个非零直和因子。
\end{solution}

\begin{exercise}\label{b2:ex:11-large-corner}
设 \(k\) 为域，\(A=M_3(k)\)，\(E_{ij}\) 为矩阵单位，\(e=E_{11}+E_{22}\)、\(f=E_{33}\)。写出四个空间 \(eAe,eAf,fAe,fAf\) 的矩阵形状与维数，并求 \(\operatorname{End}_A(Ae)\)。
\end{exercise}
\begin{solution}
按照前两个坐标与第三个坐标分块，四个空间依次是左上 \(2\times2\) 块、右上 \(2\times1\) 块、左下 \(1\times2\) 块、右下 \(1\times1\) 块，其他位置均为零；维数依次为 \(4,2,2,1\)，总和为九。角代数 \(eAe\cong M_2(k)\)，其单位是 \(e\)。由\cref{b2:prop:corner-endomorphism}，
\(\operatorname{End}_A(Ae)\cong M_2(k)^{\mathrm{op}}\)，所以其维数为四。这个自同态作用为右乘左上块，与 \(Ae\) 的六维底向量空间的全部自同态明显不同。
\end{solution}

\begin{exercise}\label{b2:ex:11-corner-hom}
设 \(k\) 为域，\(A=M_3(k)\)，\(E_{ij}\) 为矩阵单位，\(e=E_{11}+E_{22}\)、\(f=E_{33}\)。取
\[
b=E_{13}+E_{23}\in eAf,
\qquad c=E_{31}-E_{32}\in fAe.
\]
令 \(h_b:Ae\to Af\)、\(h_c:Af\to Ae\) 为对应的右乘同态。计算两个复合，证明其中一个为零，另一个是非零平方零自同态，并求后者的像与核的维数。复合仍按“先右边、后左边”的次序书写。
\end{exercise}
\begin{solution}
矩阵单位相乘给出
\[
bc=E_{11}-E_{12}+E_{21}-E_{22}\ne0,
\qquad cb=E_{33}-E_{33}=0.
\]
由\cref{b2:prop:corner-endomorphism}，\(h_c\circ h_b\) 是 \(Ae\) 上的右乘 \(bc\)，\(h_b\circ h_c\) 是 \(Af\) 上的右乘 \(cb\)，故后者为零。前者在 \(e\) 上取值为 \(bc\ne0\)，而 \((bc)^2=b(cb)c=0\)，所以非零且平方为零。若 \(X\in Ae\) 的前两列为 \(u,v\in k^3\)，则
\[
Xbc=\begin{pmatrix}u+v&-(u+v)&0\end{pmatrix}.
\]
其像由任意 \(w=u+v\in k^3\) 决定，维数为三；其核由 \(v=-u\) 决定，维数也为三，且像正好等于核。所有计算在特征二中仍成立。这里先作用 \(h_b\) 再作用 \(h_c\)，对应的是角块乘积 \(bc\)，与复合符号 \(h_c\circ h_b\) 的书写次序相反。
\end{solution}

\begin{exercise}\label{b2:ex:11-commuting-projections}
设 \(R\) 为含幺环，\(M\) 为幺性左 \(R\) 模，\(p,q\in\operatorname{End}_R(M)\) 都幂等且相互交换。证明四个算子
\(pq,p(1-q),(1-p)q,(1-p)(1-q)\) 是两两正交的幂等元且和为恒等，从而给出 \(M\) 的四项直和分解。给出不交换时其中某项不再幂等的例子。
\end{exercise}
\begin{solution}
因 \(p,q\) 交换，\(p,1-p\) 与 \(q,1-q\) 全部互相交换，故各乘积平方等于自身。两个不同乘积至少在一处分别含 \(p,1-p\) 或 \(q,1-q\)，相乘因此为零；展开 \(\bigl(p+(1-p)\bigr)\bigl(q+(1-q)\bigr)\) 得四项之和为恒等。由\cref{b2:prop:idempotent-module-decomposition}，四个像组成内直和，允许某些像为零。

在 \(k^2\) 上取 \(p=\begin{pmatrix}1&0\\0&0\end{pmatrix}\)、\(q=\begin{pmatrix}1&1\\0&0\end{pmatrix}\)。它们都幂等，但 \(pq=q\)、\(qp=p\)，不交换。此时 \(p(1-q)=\begin{pmatrix}0&-1\\0&0\end{pmatrix}\) 非零而平方为零，所以不是幂等元。原分解不能不加交换条件照搬。
\end{solution}

\begin{exercise}\label{b2:ex:11-cyclic-operator-centralizer}
设 \(k\) 为域，\(V\) 为 \(n\ge1\) 维 \(k\) 向量空间，\(T\in\operatorname{End}_k(V)\) 有循环向量 \(v\)，即 \(v,Tv,\ldots,T^{n-1}v\) 是 \(V\) 的基。所有中心化子均在 \(\operatorname{End}_k(V)\) 中取。证明 \(k[T]'=k[T]\)，从而 \(k[T]''=k[T]\)。写出三阶幂零 Jordan 块在 \(M_3(k)\) 中的中心化子。
\end{exercise}
\begin{solution}
若 \(S\) 与 \(T\) 交换，将 \(S(v)\) 按循环基展开，唯一写成 \(p(T)v\)，其中 \(\deg p<n\)。于是
\(S(T^jv)=T^jS(v)=p(T)T^jv\)，所以 \(S=p(T)\)。反过来，\(T\) 的任意多项式都与 \(T\) 交换，得到第一个等式；再取一次中心化子便得第二式。对标准上三角的 \(J_3\)，中心化子中的元素恰为
\[
aI+bJ_3+cJ_3^2
=\begin{pmatrix}a&b&c\\0&a&b\\0&0&a\end{pmatrix},\qquad a,b,c\in k.
\]
它包含非零幂零元，说明双中心化子等式本身不意味着作用代数没有幂零元。
\end{solution}

\begin{exercise}\label{b2:ex:11-repeated-matrix-blocks}
设 \(k\) 为域、\(n\ge1\)，在 \(V=k^n\oplus k^n\) 上令 \(\mathcal A=\bigl\{\operatorname{diag}(B,B):B\in M_n(k)\bigr\}\)。在 \(\operatorname{End}_k(V)\) 中求 \(\mathcal A'\) 与 \(\mathcal A''\)，并给出两者维数。
\end{exercise}
\begin{solution}
把算子写成四个 \(n\times n\) 块 \(T_{ij}\)。与每个 \(\operatorname{diag}(B,B)\) 交换，等价于每个 \(T_{ij}\) 都与全部 \(B\in M_n(k)\) 交换。逐个与矩阵单位比较，得每块都是标量矩阵，因此中心化子中的元素恰为
\[
T=\begin{pmatrix}aI_n&bI_n\\cI_n&dI_n\end{pmatrix},\qquad a,b,c,d\in k.
\]
与其中的两个块投影交换使双中心化子为块对角，再与交换两个副本的块矩阵交换使两对角块相等，故 \(\mathcal A''=\mathcal A\)。维数分别为四与 \(n^2\)。这个结果也符合\cref{b2:prop:tensor-factor-centralizers}在 \(k^n\otimes k^2\) 上的两个独立因子作用。
\end{solution}

\begin{exercise}\label{b2:ex:11-two-eigenvalue-blocks}
设 \(k\) 为域，\(\lambda,\mu\in k\)、\(\lambda\ne\mu\)，\(r,s>0\)，\(T=\operatorname{diag}(\lambda I_r,\mu I_s)\) 作用在 \(k^{r+s}\) 上。在 \(M_{r+s}(k)\) 中计算 \(k[T]'\) 和 \(k[T]''\)，并说明在 \(r>1\) 或 \(s>1\) 时，这与循环算子的例子有何不同。
\end{exercise}
\begin{solution}
分块交换方程使两个非对角块分别满足 \((\lambda-\mu)X=0\)、\((\mu-\lambda)Y=0\)，故为零；对角块任意，所以
\[
k[T]'=M_r(k)\times M_s(k)
\]
以块对角方式作用。其中心化子是每块各一个标量，故 \(k[T]''=kI_r\times kI_s\)。线性多项式可以在不同的 \(\lambda,\mu\) 处取任意两个给定值，所以后一代数恰为 \(k[T]\)。当有重数大于一时，第一次中心化子已经严格大于 \(k[T]\)；不过第二次中心化仍能恢复 \(k[T]\)，并不要求第一次中心化子等于自身。
\end{solution}

\begin{exercise}\label{b2:ex:11-tensor-flip-centralizer}
设 \(k\) 为特征不为二的域、\(V=k^2\)，令 \(\tau(v\otimes w)=w\otimes v\) 作用在 \(V\otimes_kV\) 上。在 \(\operatorname{End}_k(V\otimes_kV)\) 中求 \(k[\tau]'\) 与 \(k[\tau]''\)，并计算维数。把两个特征子空间与对称幂、外幂联系起来。
\end{exercise}
\begin{solution}
\(\tau\) 的 \(1\) 特征子空间有基
\(e_1\otimes e_1,e_2\otimes e_2,e_1\otimes e_2+e_2\otimes e_1\)，维数三；\(-1\) 特征子空间由 \(e_1\otimes e_2-e_2\otimes e_1\) 生成，维数一。特征不为二保证这些向量组成四维空间的一组基。与 \(\tau\) 交换恰好要求分别保持这两个特征子空间，因此中心化子为 \(M_3(k)\times k\)，维数十。再取中心化子，得到两个特征空间上的独立标量作用，即
\[
k\frac{1+\tau}{2}\oplus k\frac{1-\tau}{2}=kI\oplus k\tau,
\]
维数二。由\cref{b2:prop:symmetrization-characteristic}，\(\operatorname{Sym}^2_kV\) 通过 \(vw\mapsto(v\otimes w+w\otimes v)/2\) 与正特征子空间自然同构。映射 \(v\wedge w\mapsto(v\otimes w-w\otimes v)/2\) 由交替双线性关系良定义，取商到 \(\bigwedge^2_kV\) 又把它送回 \(v\wedge w\)，故它将外幂同构到负特征子空间。因而这里的两个块分别记录对称张量和交替张量上的独立算子。
\end{solution}
